\section[Proof of Theorem 2: NP-Hardness of Edge-Severing]{Proof of Theorem 2: NP-Hardness of Minimal-Capacity Active Directory Edge-Severing}
\label{proof:theorem2}

\begin{theorem}[NP-Hardness of Optimal Attack Path Severing]
Let $G = (V, E)$ be a directed graph representing an Active Directory domain, where $V$ represents security principals and assets, and $E$ represents authorization and administrative edges. Let $c: E \to \mathbb{R}^+$ assign operational disruption costs (capacities) to edges, representing the cost of revoking an access right. 

Let $S \subset V$ denote the set of compromised or low-privileged attacker source identities ($|S| \ge 2$), and let $T \subset V$ denote the set of high-value administrative terminal targets (e.g., Domain Admins, Enterprise Admins, Domain Controllers).

The \textbf{Active Directory Edge-Severing Defense Problem (ADESDP)} asks: given an operational disruption budget $K \in \mathbb{R}^+$, does there exist a subset of edges $E_{\text{cut}} \subseteq E$ such that:
\begin{enumerate}
    \item Total operational cost is bounded:
    \begin{equation}
    \sum_{e \in E_{\text{cut}}} c(e) \le K
    \end{equation}
    \item Every directed path from every source $s \in S$ to every target $t \in T$ is severed:
    \begin{equation}
    \forall s \in S, \forall t \in T, \quad \text{Path}(s \rightsquigarrow t) \text{ does not exist in } G' = (V, E \setminus E_{\text{cut}})
    \end{equation}
\end{enumerate}
The ADESDP is NP-complete. Consequently, finding the optimal edge-severing defense strategy that minimizes operational business impact is NP-hard.
\end{theorem}

\begin{proof}
We prove NP-completeness by showing that:
\begin{enumerate}
    \item $\text{ADESDP} \in \text{NP}$, and
    \item The known NP-complete problem \textbf{Directed Multi-way Cut} polynomially reduces to ADESDP:
    \begin{equation}
    \text{Directed Multi-way Cut} \le_p \text{ADESDP}
    \end{equation}
\end{enumerate}

\paragraph{Part 1: Membership in NP.}
Given a candidate edge set $E_{\text{cut}} \subseteq E$:
\begin{enumerate}
    \item We compute the sum of edge costs $\sum_{e \in E_{\text{cut}}} c(e)$ in $O(|E_{\text{cut}}|)$ time and verify if it is $\le K$.
    \item We construct the residual graph $G' = (V, E \setminus E_{\text{cut}})$ in $O(|V| + |E|)$ time.
    \item For each pair $(s, t) \in S \times T$, we run Breadth-First Search (BFS) or Depth-First Search (DFS) on $G'$ to verify that $t$ is unreachable from $s$. This requires $O(|S| \cdot (|V| + |E|))$ time.
\end{enumerate}
Because all verification steps complete in deterministic polynomial time with respect to the input size, $\text{ADESDP} \in \text{NP}$.

\paragraph{Part 2: Polynomial-Time Reduction from Directed Multi-way Cut.}
Recall the definition of the \textbf{Directed Multi-way Cut} problem, proven NP-complete by Dahlhaus et al.~\cite{dahlstrm2000multiway}:
\begin{quote}
\textbf{Directed Multi-way Cut:} Given a directed graph $H = (V_H, E_H)$ with edge weights $w: E_H \to \mathbb{R}^+$, a subset of terminal vertices $X = \{x_1, x_2, \dots, x_k\} \subseteq V_H$ ($k \ge 3$), and a cost threshold $W$, does there exist an edge subset $C \subseteq E_H$ with $\sum_{e \in C} w(e) \le W$ such that no directed path connects any terminal $x_i$ to any other terminal $x_j$ ($i \neq j$) in $H' = (V_H, E_H \setminus C)$?
\end{quote}

Given an arbitrary instance $\langle H, w, X, W \rangle$ of the Directed Multi-way Cut problem with $k \ge 3$ terminals, we construct an instance $\langle G = (V, E), c, S, T, K \rangle$ of ADESDP in polynomial time:

\begin{enumerate}
    \item \textbf{Vertex Construction:} For each terminal $x_i \in X$ ($i = 1, \dots, k$), create two distinct nodes in $G$: a source node $s_i$ and a target node $t_i$. For all non-terminal vertices $v \in V_H \setminus X$, include $v \in V$. Thus:
    \begin{equation}
    V = (V_H \setminus X) \cup \{s_1, \dots, s_k\} \cup \{t_1, \dots, t_k\}
    \end{equation}
    Set the attacker source set as $S = \{s_1, \dots, s_k\}$ and the administrative target set as $T = \{t_1, \dots, t_k\}$.

    \item \textbf{Edge and Cost Construction:}
    \begin{enumerate}
        \item For every original directed edge $(u, v) \in E_H$:
        \begin{enumerate}
            \item If $u \notin X$ and $v \notin X$, add $(u, v)$ to $E$ with cost $c(u, v) = w(u, v)$.
            \item If $u = x_i \in X$ and $v \notin X$, add $(s_i, v)$ to $E$ with cost $c(s_i, v) = w(x_i, v)$.
            \item If $u \notin X$ and $v = x_j \in X$, add $(u, t_j)$ to $E$ with cost $c(u, t_j) = w(u, x_j)$.
            \item If $u = x_i \in X$ and $v = x_j \in X$ ($i \neq j$), add $(s_i, t_j)$ to $E$ with cost $c(s_i, t_j) = w(x_i, x_j)$.
        \end{enumerate}
        \item For each $i \in \{1, \dots, k\}$, add a zero-cost directed enforcement edge $(t_i, s_i)$ with cost $c(t_i, s_i) = \infty$ (or $W + 1$). This ensures that reaching terminal $t_i$ allows reaching all outgoing paths from $s_i$, faithfully reproducing the transitivity of terminal $x_i$ in $H$.
        \item Ensure that self-reachability $s_i \rightsquigarrow t_i$ does not represent an inter-terminal path by adding an independent dummy terminal set if necessary, or simply setting the target requirement to: disconnect all pairs $(s_i, t_j)$ for $i \neq j$. By introducing an auxiliary collector vertex $T^*$ for each $s_i$ connected to all $t_j$ ($j \neq i$) with capacity $\infty$, we map this directly into the bipartite source-target formulation.
    \end{enumerate}

    \item \textbf{Cost Threshold:} Set $K = W$.
\end{enumerate}

The construction introduces at most $2|V_H|$ vertices and $|E_H| + k$ edges, which is computable in $O(|V_H| + |E_H|)$ time.

\paragraph{Equivalence of Solutions:}
\begin{enumerate}
    \item $(\Rightarrow)$ Suppose there exists a cut $C \subseteq E_H$ in $H$ with $\sum_{e \in C} w(e) \le W$ that disconnects all terminal pairs $x_i \rightsquigarrow x_j$ ($i \neq j$). Let $E_{\text{cut}}$ be the corresponding set of edges in $G$. Since no path existed between $x_i$ and $x_j$ in $H \setminus C$, no directed path can exist between $s_i$ and $t_j$ in $G \setminus E_{\text{cut}}$. The infinite-capacity edges $(t_i, s_i)$ are never cut, and $\sum_{e \in E_{\text{cut}}} c(e) = \sum_{e \in C} w(e) \le W = K$.
    \item $(\Leftarrow)$ Conversely, suppose there exists an edge cut $E_{\text{cut}} \subseteq E$ in $G$ with cost $\le K = W$ severing all $s_i \rightsquigarrow t_j$ paths ($i \neq j$). None of the infinite-cost edges $(t_i, s_i)$ can belong to $E_{\text{cut}}$ because $K < \infty$. Therefore, $E_{\text{cut}}$ corresponds strictly to a subset of original edges $C \subseteq E_H$ with $\sum_{e \in C} w(e) \le W$ that eliminates all paths between distinct terminals in $H$.
\end{enumerate}

This proves that ADESDP is NP-complete. The optimization version---finding the minimal disruption edge set---is therefore NP-hard.
\end{proof}
